\chapter{Lecture \thechapter{} of 13/05/2026}

Do we have convergence? Just a few words: if we choose $\theta = 1$, set $L = \norm{A}$ ans it holds that $\tau \sigma L^2 < 1$, then $(x^k,y^k) \to (x^*,y^*)$ as $C/N$ (order 1 convergence). The proof is quite long and not very interesting.

\begin{example}[ROF model, image denoising]

    We can consider an image (black and white) as a map $g : \Omega \to [0,1]$. Imagine that in the acquisition process for $g$ an arbitrary amount of noise is added. We can recover the ''true'' image as the minimizer of the functional 

    \[ \frac{\lambda}{2} \int_\Omega \abs{u - g}^2 \, \dd x + \int_\Omega \abs{Du} \]

    where $\abs{Du}$ denotes the total variation measure. In general we seek for a minimizer $u \in \mathrm{BV}(\Omega)$, but, if $u \in W^{1,1}(\Omega)$, then $\abs{Du} = \norm{\nabla u}\dd x$

    We consider first the function $u$ to be a $W^{1,1}$ function, then, WLOG, $\Omega = [0,1] \times [0,1]$. We discretize $\Omega$ with a uniform grid in each dimension, so that $\Omega \approx (x_i,y_j)_{ij} = (ih,jh)_{ij}$ where $h$ is the discretization step. Then we define the function on the grid as a matrix of samples $u = \set{u_{ij}}_{ij}$ and the noisy function $g = \set{g_{ij}}_{ij}$. The problem then becomes 

    \[ \frac{\lambda}{2} \sum_{ij} h^2 \abs{u_{ij} - g_{ij}} + \sum_{ij} h^2 \norm{(\nabla u)_{ij}}. \]

    Let $X = \R^{M \times N}$, $Y = \left(\R^{M \times N}, \R^{M \times N}\right)$, $p \in Y \implies p = (p^1,p^2)$, and define the map $A : X \to Y$, where $Au = (\nabla u)_{ij}$ with foreward finite differences, so 

    \[ (\nabla u)_{ij}^1 = \begin{cases}
        \displaystyle\frac{u_{i+1,j} - u_{ij}}{h} \quad \text{if } i < M \\
        \\ 
        0 \quad\text{otherwise, Neumann BC}
    \end{cases}
    \]

    and 

    \[ (\nabla u)_{ij}^2 = \begin{cases}
        \displaystyle\frac{u_{i,j+1} - u_{i,j}}{h} \quad \text{if } j < N \\
        \\ 
        0 \quad\text{otherwise, Neumann BC}
    \end{cases}
    \]

    In order to apply the Chambolle-Pock algorithm, it is needed to compute the discrete adjoint operator $\nabla^*$, which could be very large to memorize in practice. Even for a $1024 \times 1024$ pixels image, it would require approximately 8 terabytes of memory to store the full dense matrix. Even when considering a sparse matrix representation, storing the explicit structure results in a significant waste of memory. 

    Instead, the adjoint operator can be implemented in a \textit{matrix-free} fashion by exploiting its analytical definition in the discrete setting. By leveraging the definition of the adjoint via the standard inner product:
    \begin{equation}
    \langle \nabla u, p \rangle = \langle u, \nabla^* p \rangle
    \end{equation}
    it can be shown that the transpose of the Forward Finite Difference (FFD) matrix used for the gradient corresponds exactly to the negative divergence operator calculated via Backward Finite Differences (BFD). 

    Therefore, rather than storing a massive matrix and performing costly matrix-vector multiplications, $\nabla^*$ can be computed on the fly using a highly efficient, pixel-by-pixel local stencil operation. This approach reduces the operator's memory footprint to zero, significantly optimizes cache reuse, and makes the algorithm ideally suited for massive parallelization on modern GPU architectures.

    Therefore we seek for the adjoint operator. Notice that 

    \begin{align*}
        \langle u,v \rangle_X &= \sum_{ij} u_{ij}v_{ij} \\
        \langle p,q \rangle_Y &= \sum_{ij} p_{ij}^1 q_{ij}^1 + p_{ij}^2 q_{ij}^2.
    \end{align*}

    The goal is to have something in the form $\sum_{ij} u_{ij} (\nabla^* p)_{ij}$. Idea: take only the terms where $u_{ij}$ appears, i.e. $-p^1_{ij},-p^2_{ij}$ and $p_{i-1,j}^1, p_{i,j-1}^2$ as when i consider the $i-1$ i get $u_{ij} - u_{i-1,j}$, therefore a term with index $(ij)$ appears. Therefore 

    \[ (\nabla^* p)_{ij} = - \left(\frac{p_{ij}^1 - p_{i-1,j}^2}{h}\right) - \left(\frac{p_{ij}^2 - p_{i,j-1}^2}{h}\right) \]

    which is the negative distcrete divergence computed with backward finite differences.

\end{example}